\documentclass[11pt]{article}

% ---------- packages ----------
\usepackage[margin=1in]{geometry}
\usepackage{amsmath,amssymb}
\usepackage[expansion=false]{microtype}
\usepackage{enumitem}
\usepackage{booktabs}
\usepackage{array}
\usepackage{xcolor}
\usepackage{listings}
\usepackage[colorlinks=true,linkcolor=blue!60!black,citecolor=blue!60!black,urlcolor=blue!60!black]{hyperref}

\lstset{
  language=Python, basicstyle=\small\ttfamily,
  keywordstyle=\color{blue!60!black}\bfseries,
  commentstyle=\color{green!40!black}\itshape,
  stringstyle=\color{red!50!black},
  showstringspaces=false, frame=single, framerule=0.2pt,
  rulecolor=\color{black!25}, xleftmargin=1em, columns=fullflexible,
  keepspaces=true, breaklines=true,
}

% ---------- macros ----------
\newcommand{\Hm}[1]{H_{#1}}
\newcommand{\F}{\mathbb{F}}
\newcommand{\Z}{\mathbb{Z}}

\title{\textbf{TopoGym: Library Overview and Environment Specifications}}
\author{Jason Carlson}
\date{August 2026}

\begin{document}
\maketitle

\noindent\textbf{Scope.} This document specifies the TopoGym
environment library in depth: the base benchmark
(\textbf{GridWorld2D}), its textured variant
(\textbf{GridWorld2D-Texture}), and its topological variant
(\textbf{GridWorld2D-Top}, non-trivial global topology via edge
identifications). Together the three variants form one benchmark, \textbf{TopoGym-v1}:
the versioned, pinned registry across all variants under the
universal action and observation spaces of
Section~\ref{sec:spaces} --- Plain, Texture, and Top are reporting
slices of a single suite, and the version pins the environment list
(future families extend to v2 without altering v1 entries). This
document is the authority on the environments themselves, and
Section~\ref{sec:code} documents the library's code interfaces with
usage examples.

% =====================================================================
\section{GridWorld2D}
\label{sec:gw}

A registry of named environments backed by one parametric generator.
Everything is generated from a frozen configuration and a seed;
nothing is hand-placed. The design goals: ground truth matches
the theory's objects exactly, so mechanism predictions and
certificates are checkable programmatically; and the interface
follows existing gym-library conventions --- stable named
environments with seed variation and a few modifiable parts --- so
the benchmark is directly usable by others.

\subsection{World model}

A \texttt{size} $\times$ \texttt{size} grid of cells, partitioned
into free and obstacle cells. The agent occupies one free cell;
actions are \{up, down, left, right\}; moving into an obstacle leaves
the agent in place. Standing conventions, enforced at generation
time:
\begin{itemize}[leftmargin=2em, itemsep=1pt]
  \item the exterior free space (the component containing the start
        cell) is 4-connected;
  \item obstacles are well-composed: no $2 \times 2$ window contains
        exactly a diagonal pair of obstacle cells (no diagonal
        pinch);
  \item every door has width one: a free cell with obstacle cells on
        two opposite sides and free cells on the other two;
  \item a chamber is an obstacle wall enclosing free interior once
        its doors are sealed; a decoy is a sealed ring enclosing no
        reachable free cell --- its interior is unvisitable by
        construction, so persistence against it is never rewarded.
\end{itemize}

\subsection{Registry: named environments (the primary interface)}
\label{sec:registry}

Following the convention of existing gym libraries, TopoGym's
canonical interface is a registry of named, pre-defined environments:
\texttt{gym.make("TopoGym/\{Family\}-\{size\}-v0", seed=$n$)}. Each
registry entry is a frozen configuration of the underlying generator;
the seed drives layout variation (placement, door positions, shape
assignment) within that configuration; and a small set of documented
parts remains modifiable as \texttt{make} kwargs (\texttt{p\_slip},
\texttt{reward\_mode}, and per-family knobs noted below). The
benchmark is the registry crossed with the published seed lists.

\begin{center}
\small
\begin{tabular}{@{}>{\raggedright\arraybackslash}p{4.6cm}
                >{\raggedright\arraybackslash}p{5.6cm}
                >{\raggedright\arraybackslash}p{3.6cm}@{}}
\toprule
\textbf{Family} & \textbf{Description} & \textbf{Sizes / variants} \\
\midrule
\texttt{Dilution-\{50,200\}} &
  one chamber, no decoys & --- \\
\texttt{Chambers2-\{50,100,200,400\}} &
  two chambers, fixed geometry; the world-scaling family & --- \\
\texttt{ChamberCount\{1,2,4,8\}-200} &
  $k$ separated chambers & kwarg: none \\
\texttt{Decoys\{0,1,2,4,8\}-50} &
  one chamber among $k_d$ sealed decoys & kwarg: \texttt{decoy\_side} \\
\texttt{Shape\{Sq,Ci,Tr,St\}-50} &
  area-matched chamber shapes & --- \\
\texttt{Nested\{1,2,3\}-50} &
  sequentially nested shells & --- \\
\texttt{GiveUp\{1,2,4\}-50} &
  door behind a dead-end corridor of the given length & --- \\
\texttt{Bottleneck\{3,6\}-100} &
  tree of rooms joined by width-1 corridors of the given length &
  kwarg: \texttt{rooms} \\
\texttt{Maze-\{50,100\}} &
  seeded perfect maze (uniform spanning tree over cells) &
  kwarg: \texttt{braid} \\
\bottomrule
\end{tabular}
\end{center}

Registry ids are stable across releases; new families extend the
registry rather than altering existing entries.

\subsection{Generator API (advanced)}
\label{sec:generator}

Under the registry sits the full parametric generator
(\texttt{topogym.generation.TopoGenConfig2D}), available to power
users and used to define registry entries. Its configuration schema:
\begin{itemize}[leftmargin=2em, itemsep=1pt]
  \item \texttt{size} --- grid side length; any integer accepted.
  \item \texttt{n\_chambers}, \texttt{n\_decoys};
        \texttt{chamber\_side}, \texttt{decoy\_side};
        \texttt{chamber\_shape}, \texttt{decoy\_shape}
        (\{square, circle, triangle, star, mixed\}, area-matched at
        equal \texttt{side} so shape is never confounded with size).
  \item \texttt{min\_sep} --- minimum pairwise Chebyshev wall
        separation, enforced at placement with an up-front packing
        feasibility check; configurations claiming the
        one-entry-per-rollout property must satisfy
        $\texttt{min\_sep} > L$ (rollout horizon), recorded per
        configuration in the manifest.
  \item \texttt{doors\_per\_chamber} (default 1);
        \texttt{door\_kind} (\texttt{open} --- a visible walkable
        width-1 doorway, the registry convention --- or \texttt{bump},
        hidden until opened by repeated bumps); door width fixed at 1
        and validated; \texttt{door\_corridor\_len} (default 0), a
        dead-end corridor outside each door controlling entry
        probability $q$.
  \item \texttt{style} --- \texttt{open} (alias \texttt{rooms}),
        \texttt{nested} (\texttt{nested\_depth},
        \texttt{shell\_spacing}), \texttt{corridor} (\texttt{rooms},
        \texttt{corridor\_len}), or \texttt{maze} (\texttt{braid});
        Section~\ref{sec:modes}.
  \item \texttt{n\_holes} / \texttt{target\_b1},
        \texttt{goal\_in\_chamber}, \texttt{seed}; \texttt{p\_slip}
        and \texttt{reward\_mode} are environment-level knobs.
\end{itemize}
A configuration serializes to the canonical string
\texttt{TG-GridWorld2D-S\{size\}-C\{c\}-D\{d\}-cs\{n\}-ds\{n\}%
-sep\{n\}-shp\{..\}-\{mode\}-slip\{p\}-seed\{n\}}, which is the
run-log key and the manifest row; registry ids are aliases for
canonical strings. Custom configurations built through this API pass
the same validity checks and ground-truth generation as registry
entries, so results on them remain comparable in kind, but only
registry entries belong to the published benchmark.

\subsection{Environment Modes}
\label{sec:modes}

\begin{itemize}[leftmargin=2em, itemsep=2pt]
  \item \texttt{open}: chambers and decoys placed in a free field by
        seeded rejection sampling subject to \texttt{min\_sep}. The
        direct realization of the theory's objects: dilution scales
        with \texttt{size}, discrimination with
        \texttt{decoy\_count}, parallel entry requires the separation
        condition. Generation: sample footprints uniformly, reject
        on overlap or separation violation, carve walls and doors,
        validate.
  \item \texttt{nested}: \texttt{depth} concentric shells around one
        innermost chamber, one door per shell, doors angularly offset
        so entry forces traversing shells in order. The sequential
        regime where the parallel-chamber separation deliberately
        fails. Bars are born sequentially: entering shell $j$ first
        exposes shell $j{+}1$'s cycle. Validity: shell spacing
        $\geq 2$; depth $\times$ spacing bounded by half the side.
  \item \texttt{corridor}: rooms joined by width-1 corridors of
        parameterized length, the room graph constrained to a tree.
        The tree constraint is load-bearing: a cyclic room graph
        would enclose wall blocks and manufacture decoy-like $\Hm{1}$
        bars, while a tree keeps free space simply connected --- zero
        environmental homology signal, only transient exploration
        pockets. The bottleneck regime; corridor length is the
        bottleneck-depth difficulty axis. Ground truth additionally
        exposes the room graph and corridor cells.
  \item \texttt{maze}: seeded perfect-maze carving (a uniform
        spanning tree over cells), the maximal-corridor-density
        cousin of \texttt{corridor} mode --- simply connected at
        \texttt{braid} $= 0$. The \texttt{braid} parameter opens
        loops with the given density; each opened loop encloses wall
        and therefore creates an incidental $\Hm{1}$ class, so
        braided mazes interpolate between the bottleneck regime and
        the discrimination regime.
\end{itemize}

\subsection{Validity checking and manifests}

The generator validates before building and rejects with a reason:
packing feasibility (footprint plus separations fits, by a
conservative area bound and then detection of rejection-sampling
failure); exterior 4-connectivity after carving; well-composedness;
door width; separation claims. The manifest tool emits a CSV of
(registry id, canonical configuration, validity, satisfied
assumptions) covering the registry --- and, for power users, any
custom configuration set. The published benchmark is the registry
manifest plus fixed seed lists --- tuning seeds and evaluation seeds,
identical environments, disjoint seeds.

\subsection{Ground truth, certification, and instrumentation}

Read-only accessors per instance: the wall mask; the chamber list
with interiors, perimeters, and door cells; the decoy list; per-cell
annotations (chamber and door membership); and exact snapshot
save/restore. The certified metadata attached to every instance
records the free space's homology, computed from the selected complex
backend (Section~\ref{sec:complex}) and cross-checked against the
analytic expectation for the placed features at generation time.
Registration-time unit tests assert the standing conventions on every
generated instance.

\paragraph{The two door conventions.} Every instance certifies its
Betti numbers under both readings of a door:
\begin{itemize}[leftmargin=2em, itemsep=1pt]
  \item \texttt{betti\_z2} (\emph{doors walkable}, the headline):
        an enclosure with a door is an enterable room, not a hole ---
        its wall is filled before computing, so $b_1$ counts only
        \emph{sealed} structure: decoys, doorless chambers, and solid
        obstacles. A \texttt{ChamberCount} world certifies
        $[1, 0, 0]$; \texttt{Decoys4} certifies $[1, 4, 0]$ (the
        four sealed decoys and nothing else); the \texttt{Top}
        surfaces, whose chambers all have doors, recover the full
        closed surface (the torus certifies $[1, 2, 1]$, with
        $H_1 = \Z^2$ and Klein's torsion $\Z \oplus \Z/2$ intact).
  \item \texttt{betti\_z2\_sealed} (\emph{doors count as walls}):
        every door is sealed, so each doored interior becomes its own
        component and every wall component --- open arcs included ---
        contributes a class. The raw free-space $b_1$ (the strict
        homotopy type of the traversable region, where even a
        C-shaped wall arc carries a class) coincides with
        $\texttt{sealed}[1]$, so no information is lost between the
        two readings.
\end{itemize}
Both are certified against independent analytic expectations at
generation time, and the accessors
\texttt{betti\_numbers\_doors\_dont\_count\_as\_walls()} /
\texttt{betti\_numbers\_doors\_count\_as\_walls()} name them
unambiguously.

\subsection{Complex backends (all variants)}
\label{sec:complex}

Every environment exposes a \texttt{complex} configuration field
selecting the geometric substrate used for free-space homology:
\begin{itemize}[leftmargin=2em, itemsep=1pt]
  \item \texttt{cubical} (default) --- the glued cubical cell complex of
        the free cells; homology is computed by GUDHI on the order
        complex of its face poset (the barycentric subdivision), which is
        robust to every edge identification of
        Section~\ref{sec:top}. This backend is the certification source
        of truth: the metadata's Betti numbers are computed here and
        cross-checked against the analytic expectation at generation
        time.
  \item \texttt{rips} --- a Vietoris--Rips complex built by GUDHI on the
        centers of the free cells at a fixed scale
        $t \in (\sqrt{2}, 2)$, under the quotient metric induced by the
        base's identification group (wrap translates and flip glide
        reflections). Orthogonal and diagonal neighbors connect, cells
        separated by a width-one wall (ambient distance 2) do not, so
        obstacles produce classes exactly as in the cubical complex; the
        backend reports dimensions 0 and 1, and the library's tests
        assert agreement with the cubical computation of the raw free
        space on every base (the strict traversable-space reading;
        its $b_1$ equals $\texttt{sealed}[1]$).
\end{itemize}
The backend governs the environment-side homology computations
(\texttt{free\_betti}, \texttt{visited\_betti},
\texttt{observed\_betti}) and is recorded in the metadata; certification
always comes from the cubical complex.

\paragraph{Determinism.} The library guarantees end-to-end determinism
up to seeds: (configuration, seed) fixes the generated layout and its
certified metadata byte-for-byte --- including everything computed
through GUDHI, since homology is exact linear algebra over
$\F_p$ with no randomness --- and (environment, reset seed, action
sequence) fixes the episode, including \texttt{p\_slip} resampling,
which draws from the episode's seeded generator. Internal iteration
orders are sorted so no result depends on interpreter hash state; a
cross-process test asserts that two fresh interpreters produce
identical layouts, metadata, and rollouts.

\subsection{Reward modes (all variants)}
\label{sec:reward}

Every benchmark variant (GridWorld2D, -Texture, -Top) supports a
\texttt{reward\_mode} configuration field. Episodes have a
pre-determined length --- $\lfloor 1.2\max(W,H)\rfloor$ steps for a
$W \times H$ world (60 on a 50-grid, 120 on a 100-grid), knowable
from the configuration alone and overridable via
\texttt{max\_steps}; short rollouts are deliberate, pairing with
lifetime coverage and teleport resets for multi-episode exploration
--- and every environment places a goal cell by
default (removable with \texttt{goal=False} for goal-free
exploration studies):
\begin{itemize}[leftmargin=2em, itemsep=1pt]
  \item \texttt{sparse} (default) --- $+1$ terminal reward on
        reaching the goal cell, placed inside a designated chamber
        (in Texture scenarios, the treasure-chest cell). Placement
        inside a chamber makes steps-to-first-reward coincide with
        steps-to-first-entry, so the reward metric stays continuous
        with the exploration metric.
  \item \texttt{none} --- no extrinsic reward; the environment is a
        pure-exploration benchmark, and event-based metrics (entry
        times, coverage) carry the evaluation.
  \item \texttt{coverage} --- $+1$ for each cell visited for the
        first time: a dense pure-exploration reward for methods that
        require a reward signal.
  \item \texttt{deceptive} --- the \texttt{sparse} target plus a
        distractor shaping gradient leading away from the target's
        door: the reward-deception regime of the novelty-search
        literature, letting reward deception be measured rather than
        only structural deception.
\end{itemize}
Ground truth exposes the reward cells and (for \texttt{deceptive})
the shaping field. Reward modes exist for community use, for training
reward-consuming baselines, and for deception studies; pure
exploration evaluations should run \texttt{none} to keep an extra
variable out of their results.

\subsection{Universal action and observation spaces}
\label{sec:spaces}

Both spaces are identical across all variants (GridWorld2D, -Texture,
-Top); an agent written against them runs on every environment in the
library unchanged.

\paragraph{Action space.} \texttt{Discrete(4)}: up, down, left,
right, always denoting \emph{screen} directions (up decreases $y$).
Moving into an obstacle leaves the agent in place. With
\texttt{p\_slip} $> 0$ the executed action is resampled uniformly
with that probability. In GridWorld2D-Top, crossing an identified
edge applies the identification map to the agent's \emph{position}
(coordinate reversal where the identification is flipped) --- the
action space and the meaning of each action are unchanged: left is
still screen-left after any crossing. (The optional egocentric
\texttt{Discrete(3)} interface instead carries a parallel-transported
frame, where a flipped seam genuinely mirrors the agent's view.) Family-specific cell mechanics (hazard
cells that reset the episode, wormhole cells that teleport;
Section~\ref{sec:tex}) are documented
with the family and never alter the action set.

\paragraph{Observation space.} The canonical observation is a fixed
vector: the agent's integer cell coordinates $(x, y)$ followed by a
texture block $t \in [0, 1]^{16}$. Slots 0--3 are reserved
library-wide for directional blocker adjacency (obstacle to the left,
right, above, below of the current cell); slots 4--15 carry
per-environment semantic features documented with each family. The
texture block is identically zero in GridWorld2D and GridWorld2D-Top;
only GridWorld2D-Texture populates it. Optional observation modes for
learned baselines --- flattened symbolic grid, or an egocentric
window of configurable radius with per-cell texture --- are available
on every variant. Snapshot save/restore is exact (deep copy), and the
environment is deterministic given (configuration, seed, action
sequence) when \texttt{p\_slip} $= 0$.

% =====================================================================
\section{GridWorld2D-Texture}
\label{sec:tex}

The textured variant (\texttt{TG-GridWorld2DTex} prefix): the same
generator with a per-cell texture channel --- semantic labels such as
blocker-adjacent (ice), door, ladder, water, and scenario-specific
others --- exposed in the observation and recorded in ground truth.
Texture is a third signal family: local like curvature, but semantic
rather than geometric, and only as trustworthy as the environment's
labeling.

\paragraph{Design intent.} Texture scenarios are built so the
taxonomy's boundary is visible: texture discriminates flagged
structure cheaply where the labeling is faithful, and fails exactly
where discrimination requires enclosure rather than adjacency. The
flavor example: arctic sailing --- most holes are ice-cube decoys
whose neighbors carry ice-adjacent texture (avoidable on sight); the
target is a treasure chest in an enclosed space reachable through a
water-only passage; texture flags ice adjacency and doorways (ice on
both sides), while whether the passage leads anywhere remains
invisible to any local feature. Scenario suites should include cells
where texture is absent or misleading, so the boundary is measured
rather than assumed.

\paragraph{Texture encoding.} Textures populate the universal
observation's texture block (Section~\ref{sec:spaces}): slots 0--3
carry directional blocker adjacency (blocker to the left, right,
above, below --- any combination), and slots 4--15 carry the
scenario's semantic features, documented per environment below.

\paragraph{Canonical environments.} Eight named environments form the
Texture slice of TopoGym-v1:
\begin{itemize}[leftmargin=2em, itemsep=3pt]
  \item \texttt{TopoGym/IceShip}. Arctic ship navigation (the agent
        is the sailboat, and \emph{hitting ice ends the episode} ---
        collisions hurt the hull). Coastal ice sheets attach to the
        north and south edges --- land masses, not floating obstacles
        --- octagonal sealed bergs float in the open water as decoys,
        and the treasure sits in a cavity deep inside the large
        eastern landmass, reachable \emph{only} through a guaranteed
        narrow width-1 channel threaded across the ice (the guarantee
        is unit-tested: blocking the channel disconnects the goal).
        Ice-adjacent water cells carry the directional adjacency
        flags --- the boat's collision-warning system --- and
        open-water cells carry the water feature only. Texture makes
        berg avoidance cheap on sight; whether the channel leads
        anywhere remains invisible to any local feature.
  \item \texttt{TopoGym/Ladders}. Platforms, ladders, and bridges;
        a gem on the top platform is the target. Rooms (platforms),
        ladders, and bridges each carry their own semantic feature;
        vertical progress is only possible on ladder cells, and
        bridges connect platforms across gaps. The textured version
        of the bottleneck regime: ladders and bridges are the
        corridors, and their texture advertises them.
  \item \texttt{TopoGym/BankRobber}. Nested rooms with the money in
        the center room --- the textured version of the nested
        regime. Door cells and hallway cells carry their own
        features, so the sequential structure is advertised locally;
        the ordering constraint (which door must come first) remains
        global.
  \item \texttt{TopoGym/DontFall}. A large world with a large
        central hole: stepping onto the drop resets the episode (the
        hazard mechanic). A few dozen much smaller huts surround the
        hole, and exactly one contains the ruby. Textures
        distinguish drop-adjacent squares, regular dirt, hut
        interiors, and hut doorways. The hazard makes local novelty
        signals actively dangerous --- the most novel direction is
        the drop --- while the huts reproduce the discrimination
        regime at scale.
  \item \texttt{TopoGym/SpaceWarp}. Four chambers and a \emph{field}
        of wormholes: 30+ paired teleporter cells scattered evenly
        across the map on a jittered lattice (kwarg
        \texttt{warp\_sep} sets the minimum spacing; never on or next
        to a doorway, so every door remains enterable). Stepping onto
        a wormhole teleports the agent to its partner (the teleport
        mechanic). The treasure chamber has no door at all: from
        outside it is indistinguishable from a sealed decoy, and it
        is enterable only through one wormhole pair running
        chamber-interior to chamber-interior. Texture flags a cell as
        a wormhole but never which one --- all wormholes are purple
        --- so wormholes are recognizable on sight while their
        destinations must be explored. The field is the point: it is
        constant noise for any gradient-follower, while methods that
        track local transition structure can model each jump once
        traversed. Reachability of every chamber under wormhole
        transitions is guaranteed and unit-tested.
  \item \texttt{TopoGym/ClownChase}. The reward-deception scenario.
        Three sealed decoy tents cluster on one side of the map with
        a troupe of clown NPCs (\texttt{n\_clowns} make-kwarg,
        default two) random-walking among them, each leashed to its
        own tent; every agent step that reduces the distance to the
        \emph{nearest} clown pays a tiny reward ($10^{-3}$) from one
        shared budget of $2.0$ that spans the agent's lifetime on the
        layout --- the distractors run out after a few thousand
        rewarding steps in total. The treasure chamber sits on the
        opposite side (the construction rejects layouts where the
        sides are not genuinely opposite). Texture distinguishes
        doors, floor, room interiors, the treasure cell, and clown
        proximity (a slot lights up when any clown is within one
        cell); the clowns are otherwise sensed only through the
        reward gradient --- which is exactly the signal that must be
        distrusted.
  \item \texttt{TopoGym/EnvironmentalIceShip}. IceShip with seasons
        and a structured landmass: three cavities --- one holding the
        treasure, two empty --- each behind its own guaranteed
        channel, plus two sealed water pockets (little enclosed
        lakes, visible but unreachable: the certified $b_0 = 3$
        counts them). Each episode is drawn sunny (summer) or snowing
        (winter), with matching palettes; the water texture's value
        encodes the season (1.0 sunny, 0.5 cold). Only the
        \emph{floating bergs} respond to the season: in winter they
        grow --- their water fringe freezes in waves at steps 20 and
        40, and being on a cell when it freezes ends the episode ---
        and in summer their rims melt away. Channels and the landmass
        never change. The certified metadata describes the
        episode-start geometry; a \texttt{season} make-kwarg forces
        either season for controlled evaluation.
  \item \texttt{TopoGym/SearchRescue}. Search and rescue in a
        collapsed concrete structure. A person is trapped in the one
        large intact chamber; around it, 160 rubble blocks on a dense
        lattice leave only width-1/2 passages --- a maze the rescuer
        must thread, with the start placed far from the chamber. Every
        block is a small $\Hm{1}$ class of its own (certified
        $b_1 = 161$); in the agent's own discovery filtration --- the
        archive --- rubble loops are born small and die quickly as
        their blocks are circumnavigated, while the chamber's
        enclosing class grows large and refuses to die. Persistence,
        not luck, finds the person. Ten explosive red barrels sit in
        the passages (flame-marked, fatal to step on; neighbors carry
        the hazard-adjacent texture); a barrel-free route to the
        person is guaranteed.
\end{itemize}

\paragraph{SpaceWarp and local evidence (semantics, not
soundness).} SpaceWarp is built to separate two statements that
coincide everywhere else in the library: ``this enclosure admits no
local entry'' and ``this enclosure cannot be entered.'' The treasure
chamber's enclosing class is genuinely permanent in the spatial
complex (its wall has no door, so the hugging cycle never bounds even
after the interior has been visited), and exhaustive local probing of
its wall truthfully reports no entry --- yet the chamber is
enterable, through a non-local transition. Methods that classify
enclosures from local evidence alone will mark it unenterable;
methods that account for non-local transitions --- e.g.\ by analyzing
the transition graph of experience, where a wormhole edge appears
once traversed --- can do better. SpaceWarp makes that semantic
boundary measurable.

\paragraph{Semantic slot assignments.} Slots are assigned
library-wide so an agent transfers between scenarios; each scenario
uses the subset relevant to it. Cell mechanics (hazard, wormhole) are
distinct cell types, visible in every observation mode.

\begin{center}
\small
\begin{tabular}{@{}llll@{}}
\toprule
\textbf{slot} & \textbf{feature} & \textbf{slot} & \textbf{feature} \\
\midrule
0--3 & blocker adjacency (L/R/up/down) & 10 & drop-adjacent \\
4 & water & 11 & plain ground \\
5 & platform & 12 & room interior \\
6 & ladder (vertical corridor) & 13 & standing on a wormhole \\
7 & bridge (horizontal corridor) & 14 & clown within one cell \\
8 & doorway & 15 & standing on the treasure \\
9 & hallway & & \\
\bottomrule
\end{tabular}
\end{center}

Scenario sizes are fixed (IceShip/EnvironmentalIceShip/Ladders/%
BankRobber/SpaceWarp: 50; DontFall/SearchRescue: 61; ClownChase: 60); layouts vary with the seed under each
scenario's structural guarantees.

% =====================================================================
\section{GridWorld2D-Top}
\label{sec:top}

The topological variant (\texttt{TG-GridWorld2DTop} prefix):
non-trivial global topology via edge identifications of the grid's
boundary. The \texttt{topology} field selects the identification:
\begin{center}
\small
\begin{tabular}{@{}lllll@{}}
\toprule
\textbf{Value} & \textbf{Identification} & \textbf{Space} &
\textbf{$\Hm{1}(\cdot\,; \Z)$} & \textbf{$\Hm{1}(\cdot\,; \F_2)$} \\
\midrule
\texttt{plane} & none (walled boundary) & disk & $0$ & $0$ \\
\texttt{cylinder} & left--right, straight & cylinder & $\Z$ & $\F_2$ \\
\texttt{mobius} & left--right, flipped & M\"obius band & $\Z$ & $\F_2$ \\
\texttt{torus} & both pairs, straight & torus & $\Z^2$ & $\F_2^2$ \\
\texttt{klein} & one straight, one flipped & Klein bottle &
  $\Z \oplus \Z/2$ & $\F_2^2$ \\
\texttt{rp2} & both pairs, flipped & real projective plane &
  $\Z/2$ & $\F_2$ \\
\bottomrule
\end{tabular}
\end{center}
Movement wraps across identified edges, with coordinate reversal
where the identification is flipped; obstacles and chambers can be
placed on top of any topology. Note the torsion: over $\Z$ the Klein
bottle and $\mathbb{RP}^2$ carry $\Z/2$ classes that integer rank
alone misses, while $\F_2$ coefficients see them --- one reason the
library computes homology over $\F_2$.

\paragraph{Why it matters.} These spaces are locally flat everywhere:
no local signal --- curvature, texture, novelty gradient, count
statistics --- distinguishes a M\"obius band from a plane at any
single cell, because every neighborhood is Euclidean. The global
structure is visible only to signals that aggregate globally, and the
ambient $\Hm{1}$ is non-trivial with no obstacles at all. This is the
regime designed to break methods that rely on local exploration
gradients, and the cleanest possible separation between local and
global signal families.

\paragraph{Canonical layout.} The Top environments share one
scenario shape: chambers placed near the corners of the large
fundamental square, exactly one holding the treasure. The square's
sides are the edges of the identification diagram (the fundamental
polygon), so the corner regions meet across the identified edges ---
on the torus all four corners are a single point of the quotient; on
the Klein bottle and \(\mathbb{RP}^2\) they meet in pairs with
orientation reversal. What appears locally as four maximally
separated chambers is globally a tight cluster reachable across the
edges; an agent that carries only local structure treats the corners
as far apart, while the quotient makes them neighbors.

\paragraph{Registry entries.} The canonical layout is registered per
topology as \texttt{TopoGym/Top\{Plane, Cylinder, Mobius, Torus,
Klein, RP2\}-50-v0}: four chambers of side 8 near the corners of the
fundamental square, exactly one holding the treasure, start near the
center. Certified $b_1$ equals the ambient rank plus the four chamber
classes (plane and $\mathbb{RP}^2$: 4; cylinder, M\"obius, torus,
Klein: 5). The identification-aware Rips backend
(Section~\ref{sec:complex}) computes homology on the quotient metric
and is tested to agree with the cubical certification on every
topology. Detection protocols separating the ambient classes from
feature classes are left to evaluation code; the metadata provides
both the ambient base facts and the composed totals.

% =====================================================================
\section{The library: code interfaces}
\label{sec:code}

TopoGym is a Gymnasium library (\texttt{pip install topogym}; MIT).
This section documents the key abstractions with usage examples; the
repository's \texttt{docs/} tree carries per-environment pages and
the internals reference.

\subsection{The registry (the primary interface)}

Importing \texttt{topogym} registers every TopoGym-v1 id. The seed
selects the layout within the frozen configuration; documented knobs
pass through \texttt{gym.make}:

\begin{lstlisting}
import gymnasium as gym
import topogym  # registers the TopoGym/* ids

env = gym.make("TopoGym/Decoys4-50-v0", seed=3, p_slip=0.05)
obs, info = env.reset(seed=0)
info["topology"]["betti_z2"]         # [1, 4, 0]  (doors walkable)
info["topology"]["betti_z2_sealed"]  # [2, 5, 0]  (doors count as walls)

from topogym import registry
registry.registry_ids()              # every pinned id
registry.canonical_string(
    registry.get_config("TopoGym/Decoys4-50-v0"), seed=3)
# 'TG-GridWorld2D-S50-C1-D4-cs8-ds8-sep2-shpSq-open-slip0-seed3'
registry.manifest(seed=0)            # validity + assumptions per entry
\end{lstlisting}

Episodes truncate after the pre-determined $1.2\max(W,H)$ steps
(60 on 50-grids, 120 on 100-grids; Section~\ref{sec:modes}); the
sparse goal pays $+1$ by default (Section~\ref{sec:reward}).
Archive-style teleport resets are opt-in:

\begin{lstlisting}
env = gym.make("TopoGym/Maze-100-v0", seed=1, teleport=True)
env.reset(seed=0)
# ... explore ...
env.reset(options={"teleport": (12, 40)})  # any visited cell
\end{lstlisting}

\subsection{The generator and the fluent spec API}

Custom configurations pass the same validity checks and certification
as registry entries:

\begin{lstlisting}
from topogym.generation import TopoGenConfig2D, generate_2d

cfg = TopoGenConfig2D(
    base="klein", size=25, n_chambers=1, n_decoys=2, n_holes=0,
    chamber_shape="circle", chamber_side=8, min_sep=3,
    door_kind="open", door_corridor_len=2,
)
layout = generate_2d(cfg, seed=42)  # deterministic + certified
layout.metadata.betti_z2

from topogym.spec import Torus
env = Torus(15).holes(3).chambers(1).compile(seed=7)
\end{lstlisting}

\subsection{Measuring discovery}

\texttt{ExplorationTracker} turns a rollout into discovery-time
persistence (essential bars are the real topology of the explored
region; finite bars are transient beliefs), and \texttt{StatsRecorder}
accumulates per-episode metric rows:

\begin{lstlisting}
from topogym.stats import StatsRecorder
from topogym.tda import ExplorationTracker

env = StatsRecorder(gym.make("TopoGym/Nested3-50-v0", seed=1))
tracker = ExplorationTracker(env)
tracker.reset(seed=0)
# ... run the policy ...
tracker.summary()      # coverage, recovery steps, bar counts
env.episodes           # return, coverage, chamber entry steps, ...
m = env.metrics()      # the standardized metric set (frozen dataclass)
m.success_rate; m.sample_efficiency; m.visitation_entropy
m.mean_regret; m.planning_efficiency   # vs env.shortest_path()
m.steps_to_coverage    # {0.5: step, ..., 1.0: step}
m.steps_to_h1_holes    # {k: step the k-th loop was discovered}
\end{lstlisting}

\paragraph{The metrics interface and its toggles.} Metrics that are
potentially expensive are \emph{off by default} and enabled per
recorder; everything cheap is always on. The toggles:

\begin{itemize}[leftmargin=2em, itemsep=1pt]
  \item \texttt{record\_steps} --- keep a per-step row (reward,
        coverage, lifetime coverage) enabling
        \texttt{coverage\_at(global\_step)}; cheap but memory-linear
        in steps.
  \item \texttt{track\_holes} --- recompute the observed region's
        homology every step to timestamp discoveries:
        \texttt{steps\_to\_h0\_holes} and
        \texttt{steps\_to\_h1\_holes} map $k$ to the global step at
        which the observed region first carried $h_i \ge k$ (runs
        GUDHI per step).
  \item \texttt{track\_curvature} --- include
        \texttt{curvature\_coverage\_below\_zero}: the fraction of
        negatively Ollivier--Ricci-curved cells (corridors, doorways,
        bottlenecks) the agent has reached. Curvature is exact
        ($\alpha = 0$, unit-expanded $W_1$ via Hungarian matching),
        computed once per layout and cached;
        \texttt{recorder.curvature\_coverage(x)} answers ``percent
        of cells with curvature $< x$ reached'' for any threshold
        regardless of the toggle, and \texttt{env.ollivier\_ricci()}
        exposes the raw per-cell field.
\end{itemize}

\begin{lstlisting}
env = StatsRecorder(gym.make("TopoGym/Decoys4-50-v0", seed=1),
                    record_steps=True, track_holes=True,
                    track_curvature=True)
# ... run episodes ...
m = env.metrics()          # frozen dataclass; m.to_dict() for logging
m.steps_to_h1_holes        # {k: global step the k-th loop was found}
m.curvature_coverage_below_zero
env.curvature_coverage(-0.2)   # any threshold, on demand
env.unwrapped.homology_stats("observed")  # HomologyStats(h0, h1, h2=None)
\end{lstlisting}

\texttt{env.homology\_stats(which)} returns per-dimension hole
counts as a \texttt{HomologyStats} value object for
\texttt{"observed"}, \texttt{"visited"}, \texttt{"certified"}, or
\texttt{"certified\_sealed"}; gridworlds carry $h_0$ and $h_1$, with
$h_2$/$h_3$ optional (\texttt{None} where the dimension does not
apply).

\paragraph{Standardized logging.} All library logging flows through
the \texttt{topogym} logger (with a \texttt{NullHandler} installed,
so it is silent until the host configures \texttt{logging}): episode
summaries at INFO, the \texttt{TOPOGYM\_DEBUG} stream at DEBUG.
\texttt{recorder.save(path)} writes the standardized run log as JSON
--- a header carrying the canonical run key
(\texttt{recorder.run\_key()}, the configuration serialized to the
canonical string of Section~\ref{sec:generator}), the library
version, and the certified topology; then the episode rows, the
metric set, and optionally the step rows. The log is a pure function
of the run --- no wall-clock timestamps --- so replaying a seed
reproduces it byte for byte.

Structure accessors ground the metrics:
\texttt{env.topology.betti\_numbers\_doors\_count\_as\_walls()} /
\texttt{\_dont\_count\_as\_walls()} return a \texttt{BettiNumbers}
value object; \texttt{env.graph()} exposes the free-cell graph as a
\texttt{networkx.Graph}; \texttt{env.shortest\_path()} the optimal
baseline; \texttt{env.bottlenecks()} the straight-through width-1
passage cells (logged at reset under \texttt{TOPOGYM\_DEBUG}).

Per-step \texttt{info} carries \texttt{coverage},
\texttt{lifetime\_coverage} (across episodes on a layout),
\texttt{chambers\_entered}, \texttt{episode\_return},
\texttt{h0\_merges}, and friends;
\texttt{env.unwrapped.chamber\_entry\_steps} gives each chamber's
first-entry step.

\subsection{Playing and rendering}

Every environment renders as \texttt{rgb\_array} (procedural
pixel-art tiles), \texttt{ansi}, or \texttt{human} (a pygame window;
\texttt{pip install topogym[play]}). The keyboard driver:

\begin{lstlisting}[language=bash]
python scripts/play.py --list
python scripts/play.py TopoGym/SpaceWarp-v0 --seed 2
# arrows move; tab reveals hidden structure; r resets;
# backspace regenerates the layout with the next seed
\end{lstlisting}

Rendering always dims cells outside the agent's current line of
sight, so what the agent can and cannot see is visible at a glance
(reveal mode shows the full map undimmed; it exists for
documentation).

\subsection{Debugging the topology live}
\label{sec:overlay}

Two environment variables turn the simulator into a topology
debugger; both are off by default and cost nothing when unset.

\paragraph{The H1 overlay.} \texttt{TOPOGYM\_OVERLAY=1} (alias
\texttt{OVERLAY\_ENABLED=1}) recomputes, on every rendered frame, the
$\Hm{1}$ classes of the \emph{strictly-visited region} --- the cells
the agent has actually stood on, which are exactly the states an
archive can restore to --- and draws each class on the grid:

\begin{itemize}[leftmargin=2em, itemsep=1pt]
  \item the \textbf{representative cycle} (yellow): the innermost
        closed loop through strictly-visited cells enclosing the
        pocket. Every cell on it is a valid archive/teleport target;
        cells merely \emph{seen} never appear. The loop is loose when
        the agent has only encircled from afar and tightens as its
        trail hugs the structure;
  \item the \textbf{rim} (green): the part of the cycle adjacent to
        seen-but-unvisited free space --- where the loop can still
        tighten. A class whose rim has gone dark is as tight as the
        walls allow; a class that dies when its pocket is walked was
        a transient belief, not a hole. Watching cycles tighten and
        transient ones collapse is the fastest way to build intuition
        for the discovery filtration of Section~\ref{sec:gw}.
\end{itemize}

A legend with the live $\Hm{1}$ count sits in the top-right corner.
\texttt{env.h1\_representatives()} returns exactly what is drawn
(one \texttt{\{cycle, rim, pocket\}} record per class, every cycle
cell strictly visited), so archive methods consume the same loops the
overlay shows --- what you see is the interface.

\begin{lstlisting}[language=bash]
# watch loops be conjectured, confirmed, and refuted as you drive:
TOPOGYM_OVERLAY=1 python scripts/play.py TopoGym/Decoys4-50-v0
\end{lstlisting}

\paragraph{The step stream.} \texttt{TOPOGYM\_DEBUG=1} streams
everything the environment computes to the console: a reset line
(layout, seed, certified Betti numbers in both door conventions,
horizon, reward mode) and one line per step with position, action,
reward, running return, termination flags, coverage, lifetime
coverage, observed fraction, known components, H0 merges, chamber
entries, doors opened, and live scenario state (season and
frozen/melted counts, clown positions and remaining budget, hazard
counts). The two variables compose:

\begin{lstlisting}[language=bash]
TOPOGYM_DEBUG=1 TOPOGYM_OVERLAY=1 \
    python scripts/play.py TopoGym/SearchRescue-v0
\end{lstlisting}

\end{document}
